\chapter{Захист кінцевих точок, моніторинг і телеметрія}
\chaptershort{Захист кінцевих точок і телеметрія}
\label{chap:ch10}

Методи виявлення (розділи~\ref{chap:ch06}--\ref{chap:ch09}) реалізуються на практиці як \term{системи захисту кінцевих точок} --- вузлів, де ЗПК зрештою виконується: робочих станцій, серверів, мобільних пристроїв. Кінцева точка --- остання й водночас найбагатша на інформацію лінія: саме тут поведінка коду спостережувана безпосередньо. Еволюція від простого антивірусу до засобів виявлення й реагування на кінцевих точках (EDR) і розширеного виявлення (XDR) відображає зміщення парадигми: від запобігання окремому файлу до безперервного моніторингу, кореляції подій і реагування в масштабі всієї інфраструктури. Розділ розглядає архітектуру захисту кінцевих точок, телеметрію й джерела подій, виявлення й реагування (EDR/XDR), профілактичні заходи (зниження поверхні атаки, контроль застосунків, найменші привілеї) та обмеження, і завершується кейсом виявлення безфайлової атаки за телеметрією. Зібрані тут дані живлять інженерію виявлення (розділ~\ref{chap:ch11}), розвідку загроз (розділ~\ref{chap:ch12}) і реагування на інциденти (розділ~\ref{chap:ch16}).

\section{Еволюція й архітектура захисту кінцевих точок}\label{sec:ch10-arch}

Класичний \term{антивірус} будувався навколо запобігання: сканування файлів за сигнатурами (розділ~\ref{chap:ch06}) під час доступу й за розкладом, блокування відомих загроз. Зростання кількості й витонченості загроз, зокрема безфайлових і living-off-the-land (розділ~\ref{chap:ch13}), виявило межі такого підходу: що не має файлу або маскується під легітимний процес, те сигнатура файлу не ловить. Відповіддю став \term{EDR} (Endpoint Detection and Response)~\cite{ch10:gartner-edr} --- безперервний запис активності кінцевої точки, виявлення підозрілих патернів поведінки й засоби реагування (ізоляція вузла, зупинення процесу, відкочування змін). \term{XDR}~\cite{ch10:xdr} розширює цю ідею, корелюючи телеметрію кінцевих точок, мережі, пошти й хмари в єдину картину.

\begin{tikzfig}
\begin{tikzpicture}[font=\scriptsize,
  b/.style={draw,rounded corners=2pt,minimum height=6mm,text width=20mm,align=center,font=\tiny},
  >={Stealth[length=1.5mm]}]
\node[b] (sens) at (0,0) {сенсори (файли, процеси, мережа, API)};
\node[b] (tel) at (0,-1.4) {телеметрія (журнал подій)};
\node[b] (det) at (3.2,-1.4) {виявлення (правила, МН, аномалії)};
\node[b] (resp) at (6.4,-1.4) {реагування (ізоляція, зупинення)};
\node[b] (cloud) at (3.2,-2.8) {кореляція (XDR): мережа, пошта, хмара};
\draw[->] (sens)--(tel); \draw[->] (tel)--(det); \draw[->] (det)--(resp);
\draw[->] (tel)--(cloud); \draw[->] (cloud)--(det);
\end{tikzpicture}
\figcap{Рис.~\thechapter.1 --- Архітектура EDR/XDR: від сенсорів кінцевої точки через телеметрію й виявлення до реагування й кореляції}
\end{tikzfig}

Сучасна кінцева точка поєднує кілька рівнів захисту (табл.~\thechapter.1): запобіжний (сигнатури, репутація, блокування), виявлювальний (поведінка, МН, аномалії), реагувальний (ізоляція, відкочування) і слідчий (запис телеметрії для розслідування). Ця багаторівневість --- пряме втілення принципу ешелонованого захисту~\cite{ch10:defense-depth}: жоден рівень не ідеальний, тому вони компенсують слабкості один одного.

Принципова новизна EDR порівняно з антивірусом --- зміщення акценту з \emph{запобігання} на \emph{видимість} і \emph{реагування}. Класичний антивірус працював за бінарною логікою «заблокувати або пропустити» в момент доступу до файлу й не зберігав історії: якщо загроза пройшла, слідів не лишалося. EDR натомість виходить із припущення, що частина загроз неодмінно обійде запобіжний рівень (припущення компрометації, assume breach), і тому безперервно записує активність, аби уможливити виявлення \emph{постфактум}, розслідування й відновлення перебігу атаки. Це змінює саму роль засобу: він стає не лише бар'єром, а й «чорною скринькою» кінцевої точки, без якої розслідування інциденту (розділ~\ref{chap:ch16}) спиралося б лише на уцілілі артефакти диска. Нормативні настанови з обробки інцидентів ЗПК~\cite{ch10:nist-edr} прямо наголошують на цінності такого запису для реагування.

Архітектурно агент EDR працює на кількох рівнях привілеїв. Компоненти в \emph{користувацькому} просторі перехоплюють виклики бібліотек і скриптові рушії; компоненти в \emph{ядрі} (драйвери-монітори, нотифікації про створення процесів і завантаження образів) дають глибшу й важчу для обходу видимість. Сучасні операційні системи надають для цього захищені інтерфейси, щоб агент отримував події надійно й водночас сам був захищений від втручання ЗПК. Ця архітектурна вимога --- наслідок того, що EDR є пріоритетною ціллю атаки: якщо зловмисник здатен вимкнути або обманути агента, уся подальша телеметрія втрачає довіру. Тому самозахист агента (цілісність, стійкість до вивантаження, виявлення спроб втручання) --- невід'ємна частина його архітектури, а не додаткова функція.

\tabhead{Таблиця~\thechapter.1 --- Рівні захисту сучасної кінцевої точки}

\begin{center}\footnotesize
\begin{tabularx}{\linewidth}{L{22mm} Y L{20mm}}
\toprule
Рівень & Функція & Приклад \\
\midrule
Запобіжний & блокування відомого & сигнатури, репутація \\
Виявлювальний & розпізнавання підозрілого & поведінка, МН \\
Реагувальний & зупинення й відновлення & ізоляція, відкочування \\
Слідчий & запис для розслідування & телеметрія EDR \\
\bottomrule
\end{tabularx}
\end{center}

\section{Телеметрія й джерела подій}\label{sec:ch10-telemetry}

Основа EDR --- \term{телеметрія}: безперервний запис подій кінцевої точки. Ключові джерела в екосистемі Windows: \emph{трасування подій} (ETW~\cite{ch10:etw}) --- вбудований механізм, що постачає події процесів, потоків, мережі, реєстру, завантаження модулів; Sysmon~\cite{ch10:sysmon} --- доповнення, що дає докладні, зручні для виявлення події (створення процесів із командним рядком і гешами, мережеві з'єднання, зміни реєстру, створення віддалених потоків --- ознака ін'єкцій, розділ~\ref{chap:ch05}); \emph{інтерфейс сканування} (AMSI~\cite{ch10:amsi}) --- точка перехоплення скриптів і коду в пам'яті \emph{після} деобфускації, що критично для виявлення безфайлових атак і обфускованих скриптів (розділ~\ref{chap:ch08}). Аналоги в інших ОС --- аудит системних викликів, фреймворки спостереження ядра.

Цінність телеметрії визначається трьома чинниками: \emph{повнотою} (чи фіксуються потрібні події), \emph{достовірністю} (чи не може ЗПК її підробити або вимкнути) і \emph{придатністю до аналізу} (структурованість, збагачення контекстом). Критична подія --- створення процесу з повним командним рядком і батьківським процесом: саме ланцюги породження процесів («офісний застосунок $\to$ інтерпретатор команд $\to$ завантажувач») є характерним слідом атак. AMSI особливо цінний, бо перехоплює код \emph{після} розгортання в пам'яті, обходячи обфускацію на рівні джерела.

\begin{tikzfig}
\begin{tikzpicture}[font=\scriptsize,
  n/.style={draw,rounded corners=2pt,minimum height=5mm,text width=26mm,align=center,font=\tiny,inner sep=2pt},
  >={Stealth[length=1.4mm]}]
\node[n] (w) at (0,0) {офісний застосунок};
\node[n] (c) at (0,-1.1) {інтерпретатор команд (обфускований рядок)};
\node[n] (d) at (0,-2.2) {завантажувач (мережеве звернення)};
\node[n] (i) at (0,-3.3) {ін'єкція в легітимний процес};
\draw[->] (w)--(c); \draw[->] (c)--(d); \draw[->] (d)--(i);
\node[n,fill=black!6] (a) at (3.8,-1.65) {AMSI: перехоплення після деобфускації};
\draw[->,dashed] (c.east) to[out=0,in=180] (a.west);
\node[n,fill=black!6] (s) at (3.8,-3.0) {Sysmon: подія створення процесу};
\draw[->,dashed] (d.east) to[out=0,in=180] (s.west);
\end{tikzpicture}
\figcap{Рис.~\thechapter.2 --- Підозрілий ланцюг породження процесів безфайлової атаки та точки спостереження телеметрії}
\end{tikzfig}

Телеметрію слід розглядати не як простий журнал, а як \emph{джерело сигналу}, з якого виявлення видобуває рідкісні ознаки серед переважно нормальної активності. Тому сучасні сенсори не лише фіксують подію, а й \emph{збагачують} її контекстом: гешем виконуваного файлу й репутацією, повним шляхом і цифровим підписом, батьківським ланцюгом, обліковим записом, геоприв'язкою мережевого з'єднання. Збагачена подія різко підвищує якість подальшого виявлення й пришвидшує розслідування, бо аналітикові не доводиться добирати контекст вручну. Водночас обсяг збагаченої телеметрії величезний, що ставить задачу балансу між глибиною запису, вартістю зберігання й швидкістю пошуку --- інженерний компроміс, який визначає практичну придатність EDR у великій інфраструктурі.

\section{Виявлення й реагування}\label{sec:ch10-edr}

На зібраній телеметрії EDR застосовує методи виявлення: \emph{правила} поведінки (розділ~\ref{chap:ch11}) --- наприклад, Sigma-правила, що описують підозрілі патерни подій; \emph{машинне навчання} (розділ~\ref{chap:ch09}) --- класифікація й виявлення аномалій; \emph{зіставлення з ATT\&CK}~\cite{mitre-attack} --- мапування подій на техніки для структурованого розуміння. Ключова перевага перед файловим скануванням --- виявлення за \emph{поведінкою} незалежно від наявності файлу, що ловить безфайлові атаки, зловживання легітимними інструментами й невідомі зразки.

Виявлена загроза запускає \term{реагування}. Автоматичне: ізоляція вузла від мережі (зупинення поширення й зв'язку з командним сервером), зупинення й карантин процесу, відкочування змін. Напівавтоматичне: сповіщення аналітика центру безпеки з контекстом для рішення. Швидкість реагування безпосередньо впливає на збитки: що менший час від виявлення до стримування, то менше встигне зробити загроза. Цей показник пов'язаний із часом присутності (dwell time, розділ~\ref{chap:ch13}): EDR скорочує обидва --- час до виявлення й час до стримування.

Сукупний час присутності загрози зручно розкласти на складові, що відповідають фазам реагування:

\begin{equation}
T_{\text{dwell}} = T_{\text{detect}} + T_{\text{triage}} + T_{\text{respond}},
\label{eq:ch10-dwell}
\end{equation}

де $T_{\text{detect}}$ --- час від компрометації до появи сигналу, $T_{\text{triage}}$ --- час на підтвердження й пріоритезацію аналітиком, $T_{\text{respond}}$ --- час на стримування. Формула~\eqref{eq:ch10-dwell} показує, що EDR скорочує dwell time не лише кращим виявленням ($T_{\text{detect}}$), а й автоматизацією триажу та реагування ($T_{\text{triage}}$, $T_{\text{respond}}$): навіть за ідеального виявлення ручний триаж і повільне реагування залишають загрозі час на дію. Саме тому автоматичне стримування й збагачення подій контекстом --- не додаткова зручність, а прямий важіль зниження збитків.

Окремо варто наголосити на ролі \emph{кореляції} в XDR. Одинична подія на кінцевій точці часто неоднозначна: запуск інтерпретатора команд сам собою законний. Але та сама подія в поєднанні з попереднім підозрілим листом, аномальним мережевим з'єднанням і наступною спробою бокового переміщення складається в однозначну картину атаки. Кореляція різнорідних джерел підвищує і повноту (видно те, що кожне джерело окремо пропускає), і точність (контекст відсікає хибні тривоги), частково пом'якшуючи проблему базової частоти саме через вимогу збігу кількох незалежних сигналів --- той самий принцип, що лежить в основі надійних правил YARA (розділ~\ref{chap:ch06}) та багатоознакового МН (розділ~\ref{chap:ch09}).

\section{Профілактичні заходи й зниження поверхні атаки}\label{sec:ch10-prevention}

Найдешевша протидія --- не допустити виконання ЗПК узагалі. \term{Зниження поверхні атаки} зменшує кількість можливих векторів: вимкнення непотрібних функцій (макроси з інтернету, застарілі протоколи), обмеження запуску скриптів, блокування типових ланцюгів експлуатації. \term{Контроль застосунків} (application control, whitelisting~\cite{ch10:appcontrol}) дозволяє виконання лише явно дозволеним програмам --- радикально ефективний проти невідомого коду, бо ЗПК просто не отримує дозволу на запуск, але вимагає дисциплінованого супроводу переліку. Принцип \term{найменших привілеїв}~\cite{ch10:least-priv} обмежує права процесів і користувачів мінімально необхідними, зменшуючи вплив скомпрометованого коду.

Профілактика має принципову перевагу перед виявленням: вона діє \emph{до} виконання й не залежить від здатності розпізнати конкретну загрозу. Контроль застосунків не потребує знати, чим саме є невідомий файл, --- достатньо, що його немає в переліку дозволених. Найменші привілеї не запобігають запуску коду, але обмежують його наслідки: завантажувач, що виконується з правами звичайного користувача, не може встановити драйвер ядра чи змінити системні області, а отже позбавлений цілого класу технік персистентності й приховування (розділ~\ref{chap:ch05}). Ця властивість робить профілактичні заходи особливо цінними проти \emph{нових} загроз, яких жоден детектор ще не знає. Їх недолік --- операційне тертя: суворий whitelisting ускладнює легітимні зміни, а жорсткі обмеження привілеїв --- адміністрування. Тому на практиці профілактику калібрують за критичністю активу: максимально сувору на серверах і спеціалізованих вузлах (розділ~\ref{chap:ch14}), поміркованішу на робочих станціях.

Окремий профілактичний рубіж --- своєчасне \emph{оновлення} та усунення вразливостей, що закриває вектори початкового доступу (експлойти), якими ЗПК потрапляє у систему. Жоден захист кінцевої точки не компенсує повністю невиправлену критичну вразливість, через яку код виконується з високими привілеями ще до того, як поведінкове виявлення встигне спрацювати. Тому гартування конфігурації, керування вразливостями й захист кінцевих точок утворюють взаємодоповнювану трійцю: перші два зменшують імовірність і наслідки компрометації, третій виявляє й стримує те, що все ж відбулося.

Ефективність цих заходів зручно оцінити через модель зниження ризику. Якщо базова ймовірність успішного ураження за один вектор становить $p$, а незалежні профілактичні заходи блокують вектори з ефективностями $e_1,\dots,e_m$, то залишкова ймовірність ураження

\begin{equation}
p_{\text{зал}} = p\cdot\prod_{j=1}^{m}(1-e_j),
\label{eq:ch10-residual}
\end{equation}

показує мультиплікативне зниження ризику за ешелонованого захисту: кожен додатковий незалежний рівень множить залишковий ризик на $(1-e_j)<1$. Формула~\eqref{eq:ch10-residual} формалізує принцип ешелонування й пояснює, чому поєднання кількох недосконалих заходів дає високу сукупну стійкість --- за умови \emph{незалежності} рівнів (спільна вразливість руйнує добуток).

\section{Обмеження й протидія ухиленню}\label{sec:ch10-limits}

Захист кінцевих точок не є абсолютним. EDR сам стає ціллю: ЗПК намагається виявити й вимкнути агента, підробити або заглушити телеметрію (зокрема обходом ETW та AMSI), діяти в прогалинах видимості. Емпіричні оцінки стійкості EDR проти технік APT~\cite{ch10:mitre-edr} показують, що проти цілеспрямованого супротивника жоден продукт не ловить усього. Крім того, EDR генерує величезний потік подій, що загострює проблему базової частоти (розділ~\ref{chap:ch03}): надмірна чутливість топить аналітиків у хибних тривогах, недостатня --- пропускає загрози. Тому EDR потребує налаштування, кваліфікованих аналітиків і процесів (розділ~\ref{chap:ch16}).

Проблема потоку подій заслуговує окремого розгляду, бо вона визначає практичну придатність EDR не менше, ніж якість алгоритмів виявлення. Велика інфраструктура генерує мільярди подій на добу, з яких переважна більшість --- нормальна активність. Виявлення мусить знайти серед них рідкісні зловмисні патерни, не затопивши аналітиків хибними тривогами. Звідси випливає кілька практичних наслідків. По-перше, правила й моделі налаштовують під конкретне середовище: те, що в одній організації є аномалією, в іншій --- буденна робота адміністраторів. По-друге, застосовують багаторівневу фільтрацію й агрегацію: дешеві правила відсівають очевидно нормальне, дорогі методи й ручний аналіз застосовують до залишку --- та сама каскадна логіка, що й у статичному аналізі (розділ~\ref{chap:ch06}). По-третє, сигнали пріоритезують за ризиком активу й достовірністю, щоб обмежений ресурс аналітиків спрямувати на найважливіше. Без цих заходів навіть технічно досконалий EDR стає марним через «втому від тривог», коли реальні загрози тонуть у шумі.

Нарешті, EDR не існує ізольовано, а є сенсором і виконавчим органом ширшої системи керування безпекою. Його телеметрія стікається в системи кореляції й аналітики (SIEM/XDR), збагачується розвідкою загроз (розділ~\ref{chap:ch12}), а рішення про реагування вбудовані в організаційні процеси (розділ~\ref{chap:ch16}) та стратегію (розділ~\ref{chap:ch17}). Технічна досконалість агента --- необхідна, але недостатня умова: без зрілих процесів, кваліфікованих людей і актуальної розвідки найкращий EDR реалізує лише частину свого потенціалу.

Фундаментальне обмеження --- залежність від \emph{видимості}: EDR бачить лише те, що фіксує телеметрія, тому прогалини в покритті (невстановлений агент, непідтримувана платформа, вимкнене джерело подій) є сліпими зонами. Звідси випливає принцип: захист кінцевих точок доповнюють мережевим моніторингом і сегментацією (щоб бачити й стримувати те, що уникло кінцевої точки), а також гартуванням конфігурації, щоб зменшити залежність від реактивного виявлення. Вітчизняні дослідження розглядають архітектуру й можливості EDR у корпоративних мережах~\cite{ch10:edr-ua} та моніторинг аномальної активності на вузлах~\cite{ch10:monitor-ua}.

\section{Кейс: виявлення безфайлової атаки за телеметрією}\label{sec:ch10-case}

Розглянемо виявлення безфайлової атаки, що не залишає файлу на диску й тому невидима файловому скануванню (захисна задача).

\emph{Крок 1. Збір телеметрії.} EDR фіксує подію: офісний документ породжує інтерпретатор команд із підозрілим, обфускованим командним рядком --- нетиповий ланцюг породження процесів.

\emph{Крок 2. Перехоплення в пам'яті.} AMSI~\cite{ch10:amsi} перехоплює скрипт \emph{після} деобфускації, розкриваючи справжній код завантажувача, попри обфускацію джерела (розділ~\ref{chap:ch08}).

\emph{Крок 3. Кореляція.} Система зіставляє події: батьківський процес (офісний застосунок), командний рядок, мережеве звернення до зовнішнього вузла, спробу ін'єкції в легітимний процес --- і мапує їх на техніки ATT\&CK, формуючи цілісну картину атаки.

\emph{Крок 4. Реагування.} За правилом спрацьовує автоматичне стримування: ізоляція вузла від мережі, зупинення ланцюга процесів, збереження артефактів пам'яті для розслідування (розділ~\ref{chap:ch16}).

\emph{Крок 5. Збагачення й поширення.} З інциденту видобувають індикатори (адреси, геші, патерни команд), передають у розвідку загроз (розділ~\ref{chap:ch12}) і правила виявлення (розділ~\ref{chap:ch11}) для захисту решти вузлів.

Кейс демонструє силу поведінкового моніторингу: безфайлова атака, невидима для файлового сканування, виявляється за ланцюгом подій і перехопленням у пам'яті, а кореляція й швидке реагування мінімізують збитки. Водночас успіх залежить від повноти телеметрії --- прогалина в будь-якій ланці стала б сліпою зоною.

\section*{Висновки до розділу}
\addcontentsline{toc}{section}{Висновки до розділу}

Захист кінцевих точок реалізує методи виявлення там, де ЗПК виконується, і еволюціонував від файлового антивірусу до EDR і XDR --- безперервного моніторингу, кореляції подій і реагування в масштабі інфраструктури. Сучасна кінцева точка поєднує запобіжний, виявлювальний, реагувальний і слідчий рівні (табл.~\thechapter.1) за принципом ешелонованого захисту. Основа EDR --- телеметрія (ETW, Sysmon, AMSI та аналоги), цінність якої визначають повнота, достовірність і придатність до аналізу; ключовим слідом атак є ланцюги породження процесів, а AMSI перехоплює код після деобфускації. На телеметрії застосовують правила, МН і зіставлення з ATT\&CK, виявляючи загрози за поведінкою незалежно від наявності файлу, із швидким автоматичним реагуванням. Профілактика --- зниження поверхні атаки, контроль застосунків і найменші привілеї --- дає мультиплікативне зниження ризику~\eqref{eq:ch10-residual} за незалежності рівнів. Обмеження --- EDR сам є ціллю, залежить від видимості й генерує потік подій, що загострює проблему базової частоти. Головний висновок: захист кінцевих точок --- центральна, але не самодостатня ланка; його сила у поєднанні поведінкового моніторингу, профілактики, кореляції й швидкого реагування, доповнених мережевим рівнем і кваліфікованими процесами.

\selfcheck
\begin{enumerate}
\item Чому кінцева точка є найбагатшою на інформацію лінією захисту? Як еволюціонував захист від антивірусу до EDR/XDR?
\item У чому межі класичного файлового антивірусу? Які загрози він принципово не ловить?
\item Опишіть рівні захисту сучасної кінцевої точки (табл.~\thechapter.1). Як вони втілюють ешелонований захист?
\item Що таке телеметрія? Назвіть ключові джерела подій (ETW, Sysmon, AMSI) та їхню роль.
\item Чому ланцюг породження процесів є характерним слідом атак? Чим особливо цінний AMSI?
\item Які методи виявлення застосовує EDR на телеметрії? У чому перевага виявлення за поведінкою перед файловим скануванням?
\item Опишіть види реагування EDR. Як швидкість реагування пов'язана з часом присутності (dwell time)?
\item Що таке зниження поверхні атаки, контроль застосунків і найменші привілеї? Чому whitelisting радикально ефективний?
\item Поясніть формулу залишкового ризику~\eqref{eq:ch10-residual}. Чому незалежність рівнів критична для ешелонування?
\item Якими способами ЗПК протидіє EDR? Чому жоден продукт не ловить усього проти APT?
\item Як проблема базової частоти проявляється в потоці подій EDR? Що потрібно для її подолання?
\item Поясніть залежність EDR від видимості. Чому захист кінцевих точок доповнюють мережевим рівнем?
\item Опишіть виявлення безфайлової атаки за телеметрією (кейс підрозділу~\ref{sec:ch10-case}). Яку роль відіграють AMSI та кореляція?
\item Дослідницьке: оцініть покриття технік ATT\&CK телеметрією обраної конфігурації Sysmon і запропонуйте доповнення для усунення сліпих зон.
\item Дослідницьке: дослідіть вплив налаштування чутливості правил EDR на баланс повноти й частоти хибних тривог на реалістичному потоці подій.
\end{enumerate}
